%%%%%%%%%%%%%%%%%%%%%%%%%%%%%%%%%%%%%%
%% Summary of Analysis Tools in Data Structues & Algorithms in Java.
%% =================================================
%% filename: 	analysis-tools.tex
%% created:		20.04.2016
%% -----------------------------------------------------------------
%% (C) 2016 by 	Nikola Stanković
%%				University of Applied Sciences Rapperswil
%%				Switzerland
%%%%%%%%%%%%%%%%%%%%%%%%%%%%%%%%%%%%%%

\documentclass[a4paper,12pt]{scrartcl}

%% Load packages
%% ========
\usepackage[utf8]{inputenc}	% this is needed for umlauts
\usepackage[ngerman]{babel} % this is needed for umlauts
\usepackage[T1]{fontenc}    % this is needed for correct output of umlauts in pdf
\usepackage{lmodern}
\usepackage[]{setspace}
\usepackage{amsmath}
\usepackage{amssymb}
\usepackage{geometry}
\usepackage{graphicx}
\usepackage[colorlinks=true]{hyperref}		% this is needed for the hyper references
\usepackage{nicefrac}		% this is needed for the skewed fractions

\onehalfspacing

%% Glossary
\usepackage[xindy]{glossaries} 
\newglossaryentry{algorithmus}
{
	name={Algorithmus},
	description={ist eine Schritt-für-Schritt Anleitung für die Ausführung einer Aufgabe, bzw. Lösung einer Problemstellung, in einer begrenzten Dauer},
	plural={Algorithmen}	
}

\newglossaryentry{datenstruktur}
{
	name={Datenstruktur},
	description={ist eine systematische Art und Weise, um Daten zu organisieren und darauf zuzugreifen},
	plural={Datenstrukturen}	
}
\makeglossaries

%% For drawing the mindmap
\usepackage{tikz}
\usetikzlibrary{mindmap,trees}

\geometry{papersize={210mm,297mm},total={160mm,240mm},top=31mm,bindingoffset=15mm}

\setlength{\headheight}{20pt}

%% Own variables
%% ========
\newcommand{\documenttitel}{Analyse Tools}
\newcommand{\documentsubtitel}{Algorithmen und Datenstrukturen 1}
\newcommand{\documentauthors}{Nikola Stanković}
\newcommand\HRule{\noindent\rule{\linewidth}{2pt}}

%% PDF meta informations
%% ========
\hypersetup{
	pdfauthor={\documentauthors},%
	pdftitle={\documenttitel},%
	pdfsubject={\documentsubtitel},%
	pdfkeywords={AD1, Analysis-Tools, Big-Oh Notation, HSR},%
	pdfproducer={LaTeX},%
	pdfcreator={pdfLaTeX}
	linktoc=all,
    linkcolor=blue
}

%% For displaying Java code with syntax highlighting
\usepackage{listings}
\usepackage{color}

\definecolor{dkgreen}{rgb}{0,0.6,0}
\definecolor{gray}{rgb}{0.5,0.5,0.5}
\definecolor{mauve}{rgb}{0.58,0,0.82}

\lstset{frame=tb,
  language=Java,
  aboveskip=3mm,
  belowskip=3mm,
  showstringspaces=false,
  columns=flexible,
  basicstyle={\small\ttfamily},
  numbers=none,
  numberstyle=\tiny\color{gray},
  keywordstyle=\color{blue},
  commentstyle=\color{dkgreen},
  stringstyle=\color{mauve},
  breaklines=true,
  breakatwhitespace=true,
  tabsize=3
}

%% For header and footer
\usepackage{calc}
\usepackage{lipsum}
\usepackage{scrpage2}
\usepackage{lastpage} 
\pagestyle{scrheadings}
\clearscrheadfoot
\automark[chapter]{section}
\ohead{\headmark}
\ihead{\documenttitel, Schwerpunkt: Big-Oh Notation}
\setheadsepline[\textwidth+5pt]{0.5pt}
\ofoot[{Seite \thepage\ von \pageref*{LastPage}}]{Seite \thepage\ von \pageref*{LastPage}}

\begin{document}

\author{
  Stanković, Nikola
}

\begin{titlepage}
	\HRule
	\vspace*{5pt}
	\begin{flushright}
	{\Huge \textbf{\documenttitel, \\ Schwerpunkt: Big-Oh Notation} \\ \vspace*{12.5pt} \Large \documentsubtitel}
	\end{flushright}
	\vspace*{5pt}
	\HRule
	\begin{flushright}
	\vspace{20pt}
	\LARGE
	\documentauthors
	\end{flushright}
	
	\vfill	
	
	\begin{center}
		\includegraphics[width=0.5\textwidth]{images/HSR_Logo_CMYK.jpg}\par\vspace{1cm}
		{\large \today \par}
	\end{center}
\end{titlepage}

\tableofcontents
\newpage

%% All content of this work
\newpage
\section{Aufgabenstellung}

Als Aufgabe des Moduls \textbf{Lerntechniken} erstelle ich eine Visualisierung für ein komplexes Thema aus dem eigenen Studiengang auf elektronischen Medien.
Als Thematik habe ich mich für die Bearbeitung des Themas \textbf{Analysis Tools} mit dem Schwerpunkt \textbf{Big-Oh Notation} aus dem Modul \textbf{\glspl{algorithmus} und \glspl{datenstruktur} 1} entschieden.
\section{Übersicht der Thematik}
Das folgende Mindmap \emph{Abbildung \ref{fig:mindmap-thematik}} soll grob die behandelten und prüfungsrelevanten Themen visualisieren.

%% Beginn des Mindmaps
\begin{figure}[!h]
\centering
	\begin{tikzpicture}
  		\path[mindmap,concept color={rgb:red,1;green,2;blue,5},text=white]
    	node[concept] {Analysis Tools} [clockwise from=-30]
    	child[concept] { 
    		node[concept] {Empirische Analyse}
    	}
    	child[concept, thick] { 
    		node[concept] {Mathematische Funktionen} 
    	}
    	child[concept, thick] { 
    		node[concept] {Big-Oh Notation} 
    	};   
	\end{tikzpicture}
	
\caption{Mindmap über die behandelte Thematik in dieser Arbeit.}
\label{fig:mindmap-thematik}
\end{figure}

Für die Modulprüfung ist es notwendig, die \textbf{O-Abschätzung} durchführen zu können. Als gegeben im Test ist das Zeitverhalten eines Algorithmus anzusehen. In der Klausur vom Frühlingssemester 2015 machte diese Aufgabe \textbf{\nicefrac{1}{9} der Gesamtpunktzahl} aus und ist somit eines der Hauptthemen.
\newpage
\section{Einführung}
\section{Empirische Analyse}

Eine Möglichkeit die Effizienz eines \gls{algorithmus} zu ermitteln, ist es diesen \gls{algorithmus} zu implementieren und experimentell durch das Ausführen zu messen. Beispielsweise mit den folgenden Code-Fragmenten \ref{lst:code-fragment-1} und \ref{lst:code-fragment-2}.

\par \medskip

\renewcommand{\lstlistingname}{Code Fragment}% Listing -> Code Fragment

\begin{lstlisting}[captionpos=b,caption={Typische Lösung für die Zeitmessung eines Algorithmus in Java},label={lst:code-fragment-1}]
long startTime = System.currentTimeMillis();
// zu testender Algorithmus ...
long elapsedTime = System.currentTimeMillis() - startTime;   
\end{lstlisting}

\par \medskip

\begin{lstlisting}[captionpos=b,caption={Zeitmessung eines Algorithmus in Java mit nanoTime()},label={lst:code-fragment-2}]
long startTime = System.nanoTime();
// zu testender Algorithmus ...
long estimatedTime = System.nanoTime() - startTime;
\end{lstlisting}

\subsection*{Probleme bzw. Grenzen der empirischen Analyse}
Das Resultat der Messung wird von Computer zu Computer und schlimmer von Ausführung zu Ausführung sich unterscheiden. Daher ist das Vergleichen von zwei \glsps{algorithmus} äusserst schwierig. Insbesondere wird der vollständig Implementierte \gls{algorithmus} vorausgesetzt und das ausprobieren von unendlich vielen Inputgrössen benötigt unendlich viel Zeit, was wir nicht aufbringen können.
\newpage
\section{Übliche Mathematische Funktionen}
\newpage
\section{Big-Oh Notation}

%Print the glossary on a new page
\newpage
\printglossaries

\end{document}